% File: quarto/kennsluakademia-template.tex
% Pandoc/Quarto template that maps erindrekar.qmd metadata onto kennsluakademia_conf.cls,
% producing erindrekar.tex / erindrekar.pdf.
\documentclass{kennsluakademia_conf}

\conference{Ráðstefna Kennsluakademíu opinberu háskólanna}
\address{Veröld Háskóla Íslands}
\date{20 nóvember, 2026}
\keywords{AI-erindrekar, gervigreind, upplýsingaverkfræði, teymisnám, sjálfstýrt
nám}

\title{Erindrekinn vinnur, nemandinn ber ábyrgð}

\author{%
Helga Ingimundardóttir\autid{}{0000-0002-2780-3546}%
}

\affil{}{Iðnaðarverkfræði-, vélaverkfræði- og tölvunarfræðideild
(IVT), Háskóli Íslands}

% Helpers pandoc's LaTeX writer expects
\providecommand{\tightlist}{%
  \setlength{\itemsep}{0pt}\setlength{\parskip}{0pt}}
\setlength{\emergencystretch}{3em}
\makeatletter
\newsavebox\pandoc@box
\newcommand*\pandocbounded[1]{%
  \sbox\pandoc@box{#1}%
  \Gscale@div\@tempa{\textheight}{\dimexpr\ht\pandoc@box+\dp\pandoc@box\relax}%
  \Gscale@div\@tempb{\linewidth}{\wd\pandoc@box}%
  \ifdim\@tempb\p@<\@tempa\p@\let\@tempa\@tempb\fi
  \ifdim\@tempa\p@<\p@\scalebox{\@tempa}{\usebox\pandoc@box}%
  \else\usebox{\pandoc@box}%
  \fi%
}
\def\fps@figure{htbp}
\makeatother
\usepackage{longtable,booktabs,array,calc}
\newcounter{none}
\makeatletter
\@ifpackageloaded{caption}{}{\usepackage{caption}}
\AtBeginDocument{%
\ifdefined\contentsname
  \renewcommand*\contentsname{Efnisyfirlit}
\else
  \newcommand\contentsname{Efnisyfirlit}
\fi
\ifdefined\listfigurename
  \renewcommand*\listfigurename{Myndayfirlit}
\else
  \newcommand\listfigurename{Myndayfirlit}
\fi
\ifdefined\listtablename
  \renewcommand*\listtablename{Töfluyfirlit}
\else
  \newcommand\listtablename{Töfluyfirlit}
\fi
\ifdefined\figurename
  \renewcommand*\figurename{Mynd}
\else
  \newcommand\figurename{Mynd}
\fi
\ifdefined\tablename
  \renewcommand*\tablename{Tafla}
\else
  \newcommand\tablename{Tafla}
\fi
}
\@ifpackageloaded{float}{}{\usepackage{float}}
\floatstyle{ruled}
\@ifundefined{c@chapter}{\newfloat{codelisting}{h}{lop}}{\newfloat{codelisting}{h}{lop}[chapter]}
\floatname{codelisting}{Listun}
\newcommand*\listoflistings{\listof{codelisting}{Listanayfirlit}}
\makeatother
\makeatletter
\makeatother
\makeatletter
\@ifpackageloaded{caption}{}{\usepackage{caption}}
\@ifpackageloaded{subcaption}{}{\usepackage{subcaption}}
\makeatother

\begin{document}

\maketitle


\section{Inngangur}\label{inngangur}

AI-erindrekar (e. \emph{AI agents}) eru að breyta því hvernig
gervigreind er notuð í hugbúnaðarþróun, gagnagreiningu og annarri
þekkingarvinnu. Ólíkt hefðbundnu spjalli eða sjálfvirkri kóðafyllingu
geta erindrekar tekið við markmiði, lesið samhengi verkefnis, notað
verkfæri og framkvæmt fjölskrefaaðgerðir, svo sem að breyta skrám, keyra
prófanir og undirbúa \emph{pull request}. Kannanir benda til þess að
notkun slíkra kerfa fari hratt vaxandi, þótt innleiðing þeirra sé víða
enn á tilraunastigi. Í alþjóðlegri könnun McKinsey sögðu 62\% svarenda
fyrirtæki sín að minnsta kosti vera að gera tilraunir með AI-erindreka,
en fæst höfðu náð að innleiða þá í stórum stíl \citep{mckinsey2025}. Í
könnun JetBrains meðal hugbúnaðarfólks sögðust 74\% nota sérhæfð
AI-verkfæri við þróun, þar á meðal kóðaaðstoð og erindreka
\citep{jetbrains2026}.

Þessi þróun kallar á nýja faglega færni. Það nægir ekki lengur að kunna
að fá gervigreind til að framleiða kóða eða greiningu. Notandinn þarf
jafnframt að geta afmarkað verkefnið, stýrt vinnuferlinu, metið og
sannreynt niðurstöður og borið ábyrgð á lokaafurðinni.

Haustið 2026 var notkun erindreka fléttuð markvisst inn í IÐN302G
Upplýsingaverkfræði, skyldunámskeið á öðru ári í grunnnámi í
iðnaðarverkfræði við Háskóla Íslands. Námskeiðið byggir upp tölvu- og
gagnalæsi með áherslu á faglegt vinnuferli: að móta spurningar, afla og
vinna gögn, halda utan um vinnuna í GitHub, gera greiningu
endurtakanlega og miðla niðurstöðum með skýrum hætti
\citep{ingimundardottir2026endurhonnun}.

Notkun erindreka er hluti af þessari víðari endurhönnun námskeiðsins.
Markmiðið er ekki aðeins að kenna nemendum að nota ný AI-verkfæri heldur
að kenna þeim að beita þeim sem hluta af faglegu vinnuferli: að afmarka
verkefni, fela erindreka viðeigandi verk, rýna afurðina og bera ábyrgð á
niðurstöðunni. Áskorunin er því að nýta þá framkvæmdargetu sem
erindrekarnir veita án þess að sjálfvirknin taki yfir þá hugsun,
ákvarðanatöku og yfirsýn sem námið á að byggja upp.

\section{Aðferð}\label{auxf0feruxf0}

IÐN302G er kennt með vendikennslu og teymisnámi (\emph{Team-Based
Learning}) og byggist upp í lotum þar sem unnið er með raunveruleg
gagnaverkefni. Nemendur taka einstaklings- og teymispróf (iRAT/tRAT),
beita aðferðum hverrar lotu á eigin gögn í GitHub (tAPP) og skila
einstaklingsígrundun (iREF).

Við þessa uppbyggingu var bætt iTA (\emph{Individual Team Assessment}).
Í upphafi nýrrar lotu er verkefnastjóri fyrri lotu spurður óundirbúinn
um kóða, ákvarðanir og vinnulag teymisins. Markmiðið er að kanna hvort
nemandinn geti útskýrt og rökstutt það sem teymið skilaði, óháð því
hvort einstök lausn var unnin af honum sjálfum, samnemanda eða með
aðstoð erindreka.

Erindrekar voru kynntir á fyrstu kennsluviku og síðan notaðir í öllum
lotum námskeiðsins. Nemendur völdu sjálfir kerfi og notuðu einkum Codex
eða Claude Code, auk GitHub Copilot. Kennslan beindist þó ekki að
einstökum verkfærum heldur að vinnulaginu: nemandinn afmarkar markmiðið,
erindrekinn leggur til eða framkvæmir breytingu og nemandinn prófar,
rýnir og tekur afstöðu til þess hvort lausnin standist.

\begin{table*}[t]
\centering
\small
\setlength{\tabcolsep}{8pt}
\renewcommand{\arraystretch}{1.2}
\begin{tabular}{p{0.44\textwidth}p{0.48\textwidth}}
\textbf{Erindrekinn getur} & \textbf{Nemandinn þarf að geta} \\
\hline
Skrifað og breytt kóða & Útskýrt hvað kóðinn gerir \\
Breytt mörgum skrám hratt & Afmarkað hvaða breytingar eru nauðsynlegar \\
Keyrt skipanir og prófanir & Metið hvort prófanirnar séu fullnægjandi \\
Samið skýrslu & Metið fullyrðingar, rök og niðurstöður \\
Unnið hratt og sjálfstætt & Stöðvað ferlið, einfaldað og breytt um stefnu \\
Búið til sannfærandi lausn & Sannreynt lausnina og tekið ábyrgð á henni \\
\end{tabular}
\caption{Verkaskipting erindrekans og nemandans: aukinni framkvæmdargetu erindrekans þarf að mæta með virkri stýringu, mati og ábyrgð nemandans.}
\end{table*}

\begin{samepage}

Nemendur nota \texttt{AGENTS} til að skilgreina verkefnareglur fyrir
erindrekann og \texttt{AGENT\_LOG} til að skrá notkun hans og mannlega
yfirferð. Í sérstakri \emph{pull request}-æfingu er grunnferlið kennt
með afmörkuðu verkefni: erindrekinn framkvæmir, en teymið þarf að lesa,
prófa og gagnrýna breytingarnar áður en þær eru samþykktar.\footnote{Dæmi
  um æfinguna og leiðbeiningar til nemenda má skoða á námskeiðsvef
  IÐN302G:
  \url{https://hi-idn.github.io/IDN302G/ai-tools/reviewing-agents.html}.}
Sama ferli er síðan endurtekið í verkefnaskilum námskeiðsins:
\textbf{afmarka, fela erindreka framkvæmd, rýna og taka ábyrgð}.

\end{samepage}

Aðgengi að verkfærunum er einnig hluti af kennsluhönnuninni. Full
erindravirkni Codex og Claude Code krefst greiddrar áskriftar. Flestir
nemendur voru þegar með slíka áskrift, sem þeir nýta einnig í öðrum
námskeiðum, en einhverjir þurftu að kaupa hana sérstaklega. Mánaðarlegur
kostnaður er um 3.000--3.500 kr. og heildarkostnaður yfir misserið því
sambærilegur við eina hefðbundna kennslubók. Til mótvægis eru engar
kennslubækur keyptar í námskeiðinu; námsefnið er opið, aðgengilegt án
endurgjalds og samið af kennara í takt við námskeiðið. Nemendur hafa
jafnframt aðgang að GitHub Copilot án endurgjalds í gegnum GitHub
Education.

Til að fá fyrstu mynd af reynslu nemenda voru opin svör úr
miðmisseriskönnun eftir sjö vikna kennslu skoðuð ásamt athugunum kennara
á vinnuferlum teyma. Gögnin eru hér notuð sem lýsandi vísbendingar um
reynslu og vinnulag, en ekki sem mæling á áhrifum kennsluhönnunarinnar.

\section{Niðurstöður}\label{niuxf0urstuxf6uxf0ur}

Fyrstu niðurstöður sýna ákveðna togstreitu. Nemendur lýsa námsefninu og
verkefnunum sem áhugaverðum og raunhæfum og nefna sérstaklega að
gagnlegt sé að kynnast erindrekum. Á sama tíma koma fram áhyggjur af
vinnuálagi, umfangi verkefna, Git-vinnuferlinu og óskýrum væntingum.
Þegar verkefni verða stór og tími knappur eykst hættan á að erindrekinn
taki yfir stærri hluta vinnunnar en nemandinn ræður við að yfirfara og
skilja.

Athuganir kennara benda í sömu átt. Þegar afmörkun og sjálfstjórn
brestur geta erindrekarnir farið of geyst: þeir framleiða meiri kóða,
flóknari lausnir og umfangsmeiri skýrslur en verkefnið krefst.
Afleiðingin getur orðið sú að nemendur missa yfirsýn yfir eigin
verkefni. Teymi sem afmarka verkefnið sjálf og fela erindrekanum smærri
og yfirfaranlegri verk virðast hins vegar halda betur utan um bæði
lausnina og eigin skilning. Þegar vandamál koma upp er jafnframt
auðveldara að aðstoða teymi sem hafa haldið verkefninu afmörkuðu en þau
sem hafa misst yfirsýn yfir umfangsmikla lausn.

Hluti af upplifðu vinnuálagi virðist þannig geta orðið til í samspili
við erindrekana sjálfa. Auðvelt er að framleiða meiri kóða, greiningu og
texta en verkefnið krefst, en námskeiðið gerir jafnframt kröfu um að
allir í teyminu geti sett sig inn í afurðina og að breytingar séu rýndar
áður en þær eru samþykktar. Eftir því sem umfangið vex verður sú rýni
tímafrekari og hætta skapast á að hún verði yfirborðskennd. Dæmi hafa
jafnframt sést um að nemendur feli öðrum erindreka að rýna breytingar.
Slík rýni getur verið gagnleg sem viðbót, en leysir ekki af hólmi
mannlega yfirferð þegar nemendur eiga sjálfir að geta útskýrt og tekið
ábyrgð á lausninni.

Þessar athuganir benda því til þess að lykilatriðið sé ekki aðeins hvort
erindreki sé notaður, heldur \textbf{hvernig vinnunni milli nemanda og
erindreka er skipt} og hversu vel verkefnið er afmarkað áður en
framkvæmd er falin erindrekanum.

\section{Umræður}\label{umruxe6uxf0ur}

Þessar fyrstu niðurstöður falla að rannsóknum sem sýna að gervigreind
getur bæði stutt nám og ýtt undir vitræna útvistun. Kerfisbundið yfirlit
og safngreining Alanazi og félaga greina frá jákvæðum áhrifum
AI-verkfæra á námsárangur í forritun, en yfirlit Farooq og félaga bendir
jafnframt á áhyggjur nemenda af of miklu trausti á verkfærin og
ónákvæmni þeirra \citep{alanazi2025, farooq2026}. Kortlagningaryfirlit
Wang og félaga dregur upp sambærilega tvíþætta mynd: spunagreind getur
stutt sjálfstjórn og þátttöku nemenda, en jafnframt veikt yfirsýn og
sjálfstætt mat ef niðurstöðum kerfanna er tekið gagnrýnislaust
\citep{wang2026}.

Kennsluhönnunaráskorunin felst því ekki aðeins í því að kenna nemendum
að nota sífellt öflugri gervigreind. Eftir því sem erindrekarnir verða
sjálfstæðari verður mikilvægara að gera mannlega stjórn, ákvarðanatöku
og ábyrgð sýnilega og metanlega. iTA, \emph{pull request}-rýni,
\texttt{AGENT\_LOG} og afmörkuð erindrekaverkefni eru tilraun til að
færa matið í þá átt.

Reynslan bendir jafnframt til þess að kenna þurfi afmörkun
erindrekaverkefna fyrr og með markvissari hætti. Áður en framkvæmd er
falin erindreka þurfa nemendur að þarfagreina verkefnið: skilgreina hvað
þarf að leysa, hvernig það tengist námsefninu og hvaða lausn svarar
matsviðmiðunum. Nemendum er frjálst að nota flóknari aðferðir en kenndar
eru í námskeiðinu, en þá þurfa þeir að skilja þær og geta rökstutt
notkun þeirra. Flóknari lausn er ekki sjálfkrafa betri lausn,
sérstaklega ef teymið missir yfirsýn yfir það sem erindrekinn hefur búið
til.

\callout{Meginspurningin færist}{%

Frá:

\begin{quote}
Hvað framleiddi teymið?
\end{quote}

Yfir í:

\begin{quote}
Hvað skilur nemandinn, hvað ákvað hann og hvað sannreyndi hann?
\end{quote}

}

Næstu skref felast því ekki endilega í meiri notkun erindreka heldur í
skýrari afmörkun verkefna, minni og yfirfaranlegri erindrekaverkefnum og
markvissari ígrundun um hvar mannleg stjórn og ábyrgð liggur í
vinnuferlinu. Nemendur þurfa ekki aðeins að læra hvernig á að fela
erindreka verkefni, heldur einnig hvenær verkefnið er orðið of stórt,
hvenær á að stöðva erindrekann og hvernig á að einfalda lausn sem hefur
vaxið úr hófi.

Eftir því sem auðveldara verður að framleiða efni verður mikilvægari
fagleg færni að afmarka hvað á að gera --- og ekki síður hvað á ekki að
gera.

\bibliographystyle{plainnat}
\bibliography{references.bib}


\end{document}
